Introducción

En julio de 1982, yo acababa de cumplir los 22 años y acababa de terminar el cuarto curso de la licenciatura de ciencias físicas, especialidad en física atómica y nuclear. Había ganado una beca para hacer prácticas de verano en la central nuclear de Cofrentes, por la época todavía en construcción. Empezaba en unos pocos días y había ido a Valencia a firmar mi contrato veraniego. Sin nada mejor que hacer aquella mañana, me entretuve en una de las librerías universitarias cercanas a Blasco Ibáñez, con el ánimo de hacerme con una buena novela con la que entretenerme durante los viajes entre Requena (donde estaba previsto que me alojara) y Cofrentes. Era una hora de trayecto de ida y otra de vuelta. No teníamos redes sociales, ni teléfonos móviles. Yo leía mucho, es cierto, pero leer no era nada tan extraño en los ochenta. El autobús que tomaría cada mañana durante los siguientes dos meses iba lleno de trabajadores de todos los oficios, desde electricistas o albañiles hasta ingenieros, pasando por estudiantes en prácticas como yo. Y casi todos llevábamos un libro a cuestas. 

El libro que encontré aquella mañana fue, nada menos, que <<El señor de los anillos>>. Eran tres volúmenes, editados por Minotauro si la memoria aún sirve y no eran baratos, pero los compré sin vacilar a cuenta de mi primer salario. Había oido hablar de Tolkien, porque alguien había contado una historia de las revueltas estudiantiles del Mayo del 68 en París, cuando en el metro aparecían pintadas que decían <<Frodo está vivo>>. Yo no sabía muy bien quién era Frodo y qué tenía que ver con la revolución que viviera, pero cualquier sombra de duda se despejó apenas leí los versos con los que arrancaba la trilogía. 

Tres anillos para los reyes elfos bajo el cielo.
Siete para los señores enanos en palacios de piedra.
Nueve para los hombres mortales condenados a morir.
Uno para el señor oscuro, sobre el trono oscuro
en la tierra de Mordor donde se extienden las sombras.

¿Por qué aquellos versos fueron suficientes para vaciarme los bolsillos y comprometerme con mil páginas de relato épico? Creo que la respuesta tiene que ver, precisamente, con el hecho de que mi cabeza estaba llena de palabras, leídas en las innumerables novelas y libros de poesía que ya había devorado por la época. 

Mi primera gran novela, si se me permite la licencia poética, fue la Biblia, en concreto el antiguo testamento, en una versión ilustrada para niños, sobre la que caía cada sábado, con seis, quizás siete años de edad, en casa de mi tía Corro. Era uno de esos libros de antes, de los que llevaban en su catálogo los vendedores ambulantes de enciclopedias. Primorosamente encuadernado, con hojas de papel satinado y láminas en color cada pocas páginas. Contaba la historia de un pueblo, protegido por un Dios poderoso, algo arbitrario y no poco cruel y los episodios eran pura acción y maravilla. Mares que se abrían al paso de las tribus y volvían a cerrarse sobre los ejércitos del Faraón, murallas que caían al resonar las trompetas, un profeta que veía una zarza ardiente que resultaba ser Yehová, que era capaz de todo tipo de acciones fabulosas e incomprensibles. Pedir a un fiel que matara a su hijo y pararle la mano en el último momento, convertir a una mujer en una estatua de sal, enviar a una banda de ángeles asesinos a que mataran a los primogénitos de los egipcios. 

Sin darme cuenta, en casa de mi tía Corro, estaba leyendo mi primera gran épica. No fue la única. Cada domingo mi padre me compraba un tebeo (las aventuras del Capitán Trueno) y un libro de la colección <<Historias Selección>>, donde descubrí a Robinson Crusoe, a Ben-Hur, a Tom Sawyer, a Nemo a Uncas, el último mohicano y a John Silver, entre otros muchos. A los quince años, cuando mi padre me regaló la Ilíada y Odisea ---esa historia abre el primer capítulo de este libro--- yo era ya una máquina de leer. A los 22, un libro de trescientas páginas caía en una semana, dos a lo sumo. Los tres volúmenes del señor de los anillos no fueron ninguna excepción. 

La trilogía de Peter Jackson llegó veinte años más tarde. Naturalmente, hice cola en el cine para verlas y naturalmente, he visto esas películas muchas veces. Las vi cuando mi hija Irene era un bebé y mi hijo Héctor aún no había nacido, las vi otra vez cuando aún eran tan niños que los abuelos nos reñían por darles sustos con los Nazgul, y otra vez cuando eran un poco más mayores, y otra vez de adolescentes y otra vez durante el Covid y otra vez recientemente. Por culpa del pesado de su padre, mis hijos habrán visto la trilogía completa no menos de cinco veces.

Y sin embargo, ninguno de los dos ha leído los libros. 

Los libros han estado siempre en casa. La versión española que compré en 1982 y la versión en inglés que compré en una librería de viejo en San Francisco hacia 1989, durante mi post-doc en California. Están muy visibles en las estanterías, provocadores y sensuales como las sirenas, con su excelente encuadernación y sus mil páginas de aventuras. EN la misma estantería están las Mil y una Noches, los poemas completos de Czeslaw Milosz y los de Raymond Carver, El Aleph de Borges, Rayuela de Cortazar, Ana Karenina de Tolstoi y los Desposeídos de Ursula K. Leguin. Todos ellos están esperando su atención, que no llega.

Irene es médico, Héctor un brillante estudiante de físicas, acabando su máster. Los dos leían de pequeños. Son de la generación de Harry Potter, pero en caso no teníamos televisión y durante su niñez, como su padre antes que ellos, los libros eran un excelente antídoto contra el aburrimiento. Ambos siguen leyendo hoy en día. Irene tiene poco tiempo, pero lo compensa con un extraño sexto sentido para seleccionar novelas interesantes. De hecho, fue ella la que me descubrió unas de las autores más interesantes contemporáneas, Elena Ferrante y leímos juntos, gozosamente, la trilogía de la amiga estupenda. Héctor acaba de leer La insoportable levedad del ser y los dos hemos pasado muchas horas hablando de la honda, casi inaudita, percepción del alma human de Kundera. 

¿Entonces? La diferencia entre mis hijos y yo no es de actitud, ni de interés. La diferencia está en los números. Aunque a ellos las redes sociales les sorprendieron relativamente tarde, no han sido inmunes a ellas, ni al asalto audiovisual sobre la lectura. En el mundo en el que yo crecí la televisión era una receta fija que ocupaba unas pocas horas del día, entre una carta de ajuste que abría a las siete o las ocho de la tarde y cerraba a las doce de la noche. No había internet y el cine era un lujo reservado a los domingos por la tarde. Yo leía antes de dormir, leía los fines de semana, me pasaba los veranos leyendo. No lo hacía por devoción intelectual o por ganar cultura. Lo hacía, a menudo, porque no había nada mejor que hacer. Pero quince años de lecturas, construyen un vocabulario y una imaginería a la que uno ya no puede sustraerse. 

Los elfos del señor de los anillos, los Nazgul, los Rohirin, las temibles batallas del abismo de Helm y el viaje final de Frodo por Mordor, se formaron en mi mente antes de ver una sola imagen. Para mi hijos y para casi todos los jóvenes del planeta, todos esos seres fantásticos, todas esas tierra misteriosas, son las que Peter Jackson imaginó. En tiempos de libros, había tantos elfos como lectores. Hoy en día, Galadriel tiene un solo rostro, Cate Blanchett, Arwen es Liv Tyler, Frodo es Elijah Wood. Son, vaya mi encendida admiración por delante, magníficas encarnaciones de los personajes de la leyenda. Pero cuando el verbo se hace carne, irremediablemente, pierde parte de su carácter divino. 

Mis hijos, como muchos jóvenes en la franja entre los veinte y los treinta años, han leído menos que sus padres, pero la lectura ha estado y sigue estando presente en sus vidas. Ahora bien, es difícil competir con el cine. El hecho de que ni Irene ni Héctor hayan leído el Señor de los Anillos (quién sabe si lo leerán alguna vez), tiene que ver mucho, creo, con la dificultad de enfrentarse a mil páginas de texto cuando ya conocen la historia, los personajes, el desenlace. Es decir, el libro ha sido aniquilado por un \emph{preemptive strike} del cine. 

Reemplazar el libro por la pantalla tiene numerosas implicaciones, no todas, ni mucho menos, perniciosas. La imaginería de Jackson en su trilogía, bien vale una misa. Blade Runner, en mi opinión, supera de largo a la novela original (Do android dreams with electric sheep?), la trilogía de El Padrino reside en una esfera del Parnaso a la que no habrían llegado jamás las novelas de Mario Puzo. Habría mucho que escribir ---y este no es el sitio--- sobre la intertextualidad del cine y la novela, sobre la miríada de formas en que se complementan y se completan. 

Pero cuando el cine ---o la serie televisiva--- reemplazan por completo a la lectura, hay, siguiendo la frase de Borges, algo que se queja. El cine, inevitablemente, es lineal ---hay sus excepciones, pero pocas, que yo sepa, han dado lugar a películas de gran difusión--- y el espectador, la mayor parte de las veces, es dócilmente pastoreado a lo largo de un educado argumento con su planteo, nudo y desenlace. Se pierde, en consecuencia, uno de los mayores placeres de la lectura que podríamos denominar <<avanzada>> y que permite el lector reconstruir o descubrir la historia ---léase Rayuela de Cortazar, para un clásico, o la reciente Historia del amor, de Nicole Krauss---. De hecho, el cine genera la obsesión por evitar el malhadado \emph{spoiler}. No queremos que nos cuenten qué pasa en la peli (no digamos ya el final) porque el argumento se convierte en el motor principal de la historia. A ese respecto, confieso sin pudor que mis amigos y familiares me temen, por mi tendencia al terrorismo gafapelis. No puedo evitarlo, en general el argumento me trae sin cuidado ---es una razón por la que nunca me han enganchado las novelas de crímenes--- y de hecho, mi técnica para comprar libros de los que sé poco o nada ---esas veces que nos llama la atención un título o un amigo nos ha recomendado una novela--- suele ser leer la primera página \emph{y la última}. Me importa como un autor acaba, igual que me importa como empieza, pero me da igual saber que Gatsby muere antes, o después de acabar la novela, porque todo el valor literario de esa obra maestra no tiene nada que ver con el detalle, obvio por otra parte desde el primer renglón, que el protagonista tiene que morir. 

Hay más cosas que se pierden en la traducción de la letra escrita a la pantalla ---traduttore traditore---. El ritmo interno de la prosa ---Fitzgerald---, el estilo ---Borge, Nabokov---, la extrañeza ---Márquez, Cortazar---, la locura ---Bulgakov--. La sensación de que Flaubert ha sacado un conejo de la chistera cuando, para demostrarte que el marido de Emma Bovary es un pobre hombre, no habla de él, sino de su sombrero. La intuición de que la realidad es un velo que esconde algo más transcendente que alimenta la lectura de Shakespeare. El simple y puro placer de repetir en voz alta unos versos perfectos. 

El ciego sol se estrella
Contra las duras aristas de las armas.

Cierto, se puede imaginar una escena en la que el desterrado Cid y su banda de fieles cabalgan bajo el sol de Castilla, pero cómo capturar la aliteración <<duras, aristas de las armas>> (repite, amable lectora, compañero lector, la línea en voz alta). O es otro:

Te acordaste de mí, cuando subías,
al silencio que sufre la serpiente
prisionera de grillos y de umbrías.

¿Cómo reducir a celuloide el infinito que despliega Lorca en tres líneas? ¿Cómo expresar la emoción con que el corazón tiembla al leer <<el silencio que sufre la serpiente>>? 

SIMPLIFICAR, UNIFORMAR PERSONAJES, TRANSMITIR IDEOLOGIA A MASAS. MISMO MENSAJE, IDÉNTICA VERDAD. 


 Escribo estas líneas a finales de agosto de 2026, 

En julio de 2026, confluyeron dos eventos que habían estado gestionándose bastante tiempo. El primero fue el estreno de La Odisea, que se ha convertido en la película de moda de e